% Options for packages loaded elsewhere
\PassOptionsToPackage{unicode}{hyperref}
\PassOptionsToPackage{hyphens}{url}
\documentclass[
 12pt,
]{ctexart}
\usepackage{xcolor}
\usepackage{amsmath,amssymb}
\setcounter{secnumdepth}{-\maxdimen} % remove section numbering
\usepackage{iftex}
\ifPDFTeX
 \usepackage[T1]{fontenc}
 \usepackage[utf8]{inputenc}
 \usepackage{textcomp} % provide euro and other symbols
\else % if luatex or xetex
 \usepackage{unicode-math} % this also loads fontspec
 \defaultfontfeatures{Scale=MatchLowercase}
 \defaultfontfeatures[\rmfamily]{Ligatures=TeX,Scale=1}
\fi
\usepackage{lmodern}
\ifPDFTeX\else
 % xetex/luatex font selection
\fi
% Use upquote if available, for straight quotes in verbatim environments
\IfFileExists{upquote.sty}{\usepackage{upquote}}{}
\IfFileExists{microtype.sty}{% use microtype if available
 \usepackage[]{microtype}
 \UseMicrotypeSet[protrusion]{basicmath} % disable protrusion for tt fonts
}{}
\makeatletter
\@ifundefined{KOMAClassName}{% if non-KOMA class
 \IfFileExists{parskip.sty}{%
 \usepackage{parskip}
 }{% else
 \setlength{\parindent}{0pt}
 \setlength{\parskip}{6pt plus 2pt minus 1pt}}
}{% if KOMA class
 \KOMAoptions{parskip=half}}
\makeatother
\usepackage{color}
\usepackage{fancyvrb}
\newcommand{\VerbBar}{|}
\newcommand{\VERB}{\Verb[commandchars=\\\{\}]}
\DefineVerbatimEnvironment{Highlighting}{Verbatim}{commandchars=\\\{\}}
% Add ',fontsize=\small' for more characters per line
\newenvironment{Shaded}{}{}
\newcommand{\AlertTok}[1]{\textcolor[rgb]{1.00,0.00,0.00}{\textbf{#1}}}
\newcommand{\AnnotationTok}[1]{\textcolor[rgb]{0.38,0.63,0.69}{\textbf{\textit{#1}}}}
\newcommand{\AttributeTok}[1]{\textcolor[rgb]{0.49,0.56,0.16}{#1}}
\newcommand{\BaseNTok}[1]{\textcolor[rgb]{0.25,0.63,0.44}{#1}}
\newcommand{\BuiltInTok}[1]{\textcolor[rgb]{0.00,0.50,0.00}{#1}}
\newcommand{\CharTok}[1]{\textcolor[rgb]{0.25,0.44,0.63}{#1}}
\newcommand{\CommentTok}[1]{\textcolor[rgb]{0.38,0.63,0.69}{\textit{#1}}}
\newcommand{\CommentVarTok}[1]{\textcolor[rgb]{0.38,0.63,0.69}{\textbf{\textit{#1}}}}
\newcommand{\ConstantTok}[1]{\textcolor[rgb]{0.53,0.00,0.00}{#1}}
\newcommand{\ControlFlowTok}[1]{\textcolor[rgb]{0.00,0.44,0.13}{\textbf{#1}}}
\newcommand{\DataTypeTok}[1]{\textcolor[rgb]{0.56,0.13,0.00}{#1}}
\newcommand{\DecValTok}[1]{\textcolor[rgb]{0.25,0.63,0.44}{#1}}
\newcommand{\DocumentationTok}[1]{\textcolor[rgb]{0.73,0.13,0.13}{\textit{#1}}}
\newcommand{\ErrorTok}[1]{\textcolor[rgb]{1.00,0.00,0.00}{\textbf{#1}}}
\newcommand{\ExtensionTok}[1]{#1}
\newcommand{\FloatTok}[1]{\textcolor[rgb]{0.25,0.63,0.44}{#1}}
\newcommand{\FunctionTok}[1]{\textcolor[rgb]{0.02,0.16,0.49}{#1}}
\newcommand{\ImportTok}[1]{\textcolor[rgb]{0.00,0.50,0.00}{\textbf{#1}}}
\newcommand{\InformationTok}[1]{\textcolor[rgb]{0.38,0.63,0.69}{\textbf{\textit{#1}}}}
\newcommand{\KeywordTok}[1]{\textcolor[rgb]{0.00,0.44,0.13}{\textbf{#1}}}
\newcommand{\NormalTok}[1]{#1}
\newcommand{\OperatorTok}[1]{\textcolor[rgb]{0.40,0.40,0.40}{#1}}
\newcommand{\OtherTok}[1]{\textcolor[rgb]{0.00,0.44,0.13}{#1}}
\newcommand{\PreprocessorTok}[1]{\textcolor[rgb]{0.74,0.48,0.00}{#1}}
\newcommand{\RegionMarkerTok}[1]{#1}
\newcommand{\SpecialCharTok}[1]{\textcolor[rgb]{0.25,0.44,0.63}{#1}}
\newcommand{\SpecialStringTok}[1]{\textcolor[rgb]{0.73,0.40,0.53}{#1}}
\newcommand{\StringTok}[1]{\textcolor[rgb]{0.25,0.44,0.63}{#1}}
\newcommand{\VariableTok}[1]{\textcolor[rgb]{0.10,0.09,0.49}{#1}}
\newcommand{\VerbatimStringTok}[1]{\textcolor[rgb]{0.25,0.44,0.63}{#1}}
\newcommand{\WarningTok}[1]{\textcolor[rgb]{0.38,0.63,0.69}{\textbf{\textit{#1}}}}
\usepackage{longtable,booktabs,array}
\usepackage{calc} % for calculating minipage widths
% Correct order of tables after \paragraph or \subparagraph
\usepackage{etoolbox}
\makeatletter
\patchcmd\longtable{\par}{\if@noskipsec\mbox{}\fi\par}{}{}
\makeatother
% Allow footnotes in longtable head/foot
\IfFileExists{footnotehyper.sty}{\usepackage{footnotehyper}}{\usepackage{footnote}}
\makesavenoteenv{longtable}
\setlength{\emergencystretch}{3em} % prevent overfull lines
\providecommand{\tightlist}{%
 \setlength{\itemsep}{0pt}\setlength{\parskip}{0pt}}
\usepackage{bookmark}
\IfFileExists{xurl.sty}{\usepackage{xurl}}{} % add URL line breaks if available
\urlstyle{same}
\hypersetup{
 hidelinks,
 pdfcreator={LaTeX via pandoc}}

\author{SCX}
\date{2026}

\begin{document}

\section{Spring
self-evolution rigorous convergence of dynamicsAnalysis}\label{spring-ux81eaux8fdbux5316ux52a8ux529bux5b66ux7684ux4e25ux5bc6ux6536ux655bux5206ux6790}

\begin{center}\rule{0.5\linewidth}{0.5pt}\end{center}

\textbf{Document Status}: Rigorous Mathematical Derivation + Honest Critique\\
\textbf{-dependence}:
\texttt{multi\_head\_spring\_and\_positional\_encoding\_analysis.md}
(CCAudit Report), \texttt{ppe\_rigorous\_derivation.md},
\texttt{hostile\_review.md}\\
\textbf{Methodology}: non-convex optimizationtheory (Robbins-Monro, Łojasiewicz, escape time) +
Online Learning (regret bound) + stochastic process (monotonicity)\\
\textbf{honestdegreeLabel}: each resultmarkNote \texttt{{[}rigorousProof{]}} /
\texttt{{[}partrigorous+Conjecture{]}} / \texttt{{[}Conjecture{]}} /
\texttt{{[}反例构造{]}}

\begin{center}\rule{0.5\linewidth}{0.5pt}\end{center}

\subsection{0.
Notation Convention and Prerequisites}\label{ux7b26ux53f7ux7ea6ux5b9aux4e0eux5148ux5907ux77e5ux8bc6}

\subsubsection{0.1 Spring
State Space}\label{spring-ux72b6ux6001ux7a7aux95f4}

Spring self-evolutiondynamics的CompleteStatusis 三元组: 

\[\boxed{\Xi_t = (\mathcal{S}_t, \theta_t, \mathcal{M}_t)}\]

where: - \(\mathcal{S}_t = \{s_1, \ldots, s_N\}\): \(N\)
 Statusatom, \(s_i \in \mathbb{R}^{d_s}\) -
\(\theta_t = \{Q_k, K_k, V_k\}_{k=1}^{K} \cup \{W^O\}\): Multi-Head
Spring allmaytrainingparameter -
\(\mathcal{M}_t = \{(s_i, \text{Cercis}_t(s_i), \text{audit}_t(s_i))\}_{i=1}^{N}\): memory bank, existence储each Statusatom的
Cercis 评separateandYajie审计result

parameterdimension: \(|\theta| = K \cdot (2d_k d_s + d_v d_s) + K d_v d_s\). 

\subsubsection{0.2
Yajie审计andDetection Margin}\label{ux96c5ux6d01ux5ba1ux8ba1ux4e0eux68c0ux6d4bux8fb9ux9645}

forStatus \(s\), \(M\) independent专家投票
\(v_m(s) = \mathbf{1}[E_m(s) \neq y]\). 

\[C(s) = \sum_{m=1}^{M} v_m(s) \sim \text{Binomial}(M, p_s)\]

\begin{itemize}
\tightlist
\item
 dry净样this (\(H_0\)): \(p_s = p_{\text{clean},s} < 0.5\)
\item
 noise样this (\(H_1\)): \(p_s = p_{\text{noisy},s} > 0.5\)
\end{itemize}

\textbf{Detection Margin} (core量): 

\[\boxed{\Delta_s(t) = p_{\text{noisy},s}(t) - p_{\text{clean},s}(t)}\]

where \(p_{\text{clean},s}(t), p_{\text{noisy},s}(t)\) isat第 \(t\)
轮进化back的专家separate歧probability. \(\Delta_s(t) > 0\) 意味着Status \(s\)
的noisemaybydetection; \(\Delta_s(t)\) the morelarge, detectionthe moremay靠. 

\subsubsection{0.3 SE-1: Robbins-Monro
随机gradient descent}\label{se-1-robbins-monro-ux968fux673aux68afux5ea6ux4e0bux964d}

Spring parameterupdate (SE-1): 

\[\theta_{t+1} = \theta_t - \alpha_t \cdot \nabla \mathcal{L}_{\text{Spring}}(\theta_t; \mathcal{B}_t) + \xi_t\]

where: - \(\mathcal{L}_{\text{Spring}}\) is Spring
self-reconstructionloss (StatusatomthroughNote意力mechanism的self-predictionerror) -
\(\mathcal{B}_t \subset \mathcal{M}_t\) isFrom memory bankmiddlesampling的small批量 -
\(\xi_t\) issmall批量gradientandfull批量gradient's差的noise项 - Robbins-Monro
步longConditions: \(\sum_{t=1}^{\infty} \alpha_t = \infty\), \(\sum_{t=1}^{\infty} \alpha_t^2 < \infty\)

\textbf{keynon-凸ness}: \(\mathcal{L}_{\text{Spring}}(\theta)\)
ishighdegreenon-凸的. Multi-Head Spring loss景观具has 排列to称ness------to任意排列
\(\pi \in S_K\): 

\[\mathcal{L}_{\text{Spring}}(\theta) = \mathcal{L}_{\text{Spring}}(\pi \cdot \theta)\]

where \(\pi \cdot \theta\) representationwill头索引按 \(\pi\) 置换. thisproduces \(K!\)
 等价的stationary point (等价类). 

\subsubsection{0.4 memory bankupdate}\label{ux8bb0ux5fc6ux5e93ux66f4ux65b0}

each轮 \(t\): 1. \textbf{evaluation}: to \(\mathcal{M}_t\) in the each Statusatom
\(s_i\), compute Cercis 评separate \(S_t(s_i)\) 2. \textbf{选择}: according to
\(S_t(s_i)\) 选取''has value''的Statusatomused for Spring training 3.
\textbf{进化}: through SE-1 morenew \(\theta_t \to \theta_{t+1}\) 4.
\textbf{审计}: heavynewcomputeYajie审计统计量 \(C(s_i)\), updateDetection Marginestimate 5.
\textbf{updatememory bank}: \(\mathcal{M}_{t+1} = \mathcal{M}_t \cup \{\text{newdiscover的Statusatom}\} \setminus \{\text{byAssessmentis purenoise的atom}\}\)

\begin{center}\rule{0.5\linewidth}{0.5pt}\end{center}

\section{Problem 1: non-凸 Robbins-Monro
convergence率}\label{ux95eeux9898-1ux975eux51f8-robbins-monro-ux6536ux655bux7387}

\subsection{1.1
ProblemExact 陈述}\label{ux95eeux9898ux7cbeux786eux9648ux8ff0}

Multi-Head Spring has \(K\) 头, loss景观has \(K!\)
 排列to称stationary point. consider带noise的gradient descent: 

\[\theta_{t+1} = \theta_t - \alpha_t \cdot (\nabla \mathcal{L}(\theta_t) + \zeta_t)\]

where
\(\mathbb{E}[\zeta_t | \mathcal{F}_t] = 0\), \(\mathbb{E}[\|\zeta_t\|^2 | \mathcal{F}_t] \leq \sigma^2\) (has 界varianceAssumptions). 

\textbf{Core Problem}: Noisy Gradient Descent iswhetherwith \(O(1/\sqrt{t})\)
rateconverges tosome stationary point的邻域? convergence率lower boundis什么? 

\subsection{1.2
loss景观的排列to称nessstructure}\label{ux635fux5931ux666fux89c2ux7684ux6392ux5217ux5bf9ux79f0ux6027ux7ed3ux6784}

\subsubsection{\texorpdfstring{1.2.1 商space构造
\texttt{{[}rigorousProof{]}}}{1.2.1 商space构造 {[}rigorousProof{]}}}\label{ux5546ux7a7aux95f4ux6784ux9020-ux4e25ux683cux8bc1ux660e}

\textbf{Definition 1.1 (排列商space)}: Definition等价relationship
\(\theta \sim \theta' \iff \exists \pi \in S_K : \theta' = \pi \cdot \theta\). 商spaceis 
\(\Theta / S_K\). 

\textbf{Proposition 1.1 (商spaceupper的gradient Lipschitz ness)}
\texttt{{[}rigorousProof{]}}: If \(\mathcal{L}\) in \(\Theta\) upperis
\(L\)-smooth滑 (i.e. \(\nabla\mathcal{L}\) is \(L\)-Lipschitz
continuous的), then诱导的商spaceupper的function
\(\bar{\mathcal{L}}([\theta]) = \mathcal{L}(\theta)\) (良Definition, becauseis 
\(\mathcal{L}\) is \(S_K\)-not变的)alsois \(L\)-smooth滑的. 

\emph{Proof}: \(\mathcal{L}\) \(S_K\)-not变nessguarantee \(\bar{\mathcal{L}}\)
is良Definition的. for \([\theta_1], [\theta_2] \in \Theta/S_K\), 选取代表元
\(\theta_1, \theta_2\) makes
\(\|\theta_1 - \theta_2\| = d([\theta_1], [\theta_2])\) (商degree量). by
\(\mathcal{L}\) in \(\Theta\) upper \(L\)-smooth滑ness: 
\[\|\nabla\bar{\mathcal{L}}([\theta_1]) - \nabla\bar{\mathcal{L}}([\theta_2])\| = \|\nabla\mathcal{L}(\theta_1) - \nabla\mathcal{L}(\theta_2)\| \leq L \cdot \|\theta_1 - \theta_2\| = L \cdot d([\theta_1], [\theta_2])\]
where第一 等号利using \(\nabla\mathcal{L}\) 
\(S_K\)-等变ness (\(\nabla\mathcal{L}(\pi \cdot \theta) = \pi \cdot \nabla\mathcal{L}(\theta)\)), thisguarantee了
\(\nabla\bar{\mathcal{L}}\) is良Definition的. \(\square\)

\textbf{Proposition 1.2 (stationary point的separate类)}
\texttt{{[}rigorousProof{]}}: \(\Theta / S_K\) in the each stationary point \([\theta^*]\)
toshould \(\Theta\) in the \(K!\) 孤立stationary point (whenstable子群
\(\text{Stab}_{S_K}(\theta^*) = \{e\}\) time), or \(K! / |\text{Stab}|\)
 stationary point (whenexistenceatto称nesstime, such assome些头坍塌tomutualsameparameter). 

\emph{Proof}: thisis群作用的轨道-stable子Theorem的直接Corollary. \(S_K\) in
\(\Theta\) upper的作用is \(\mathcal{L}\) to称群. 轨道size =
\(|S_K| / |\text{Stab}_{S_K}(\theta)|\). \(\square\)

\subsubsection{\texorpdfstring{1.2.2 saddle point逃逸Analysis
\texttt{{[}rigorousProof{]}}}{1.2.2 saddle point逃逸Analysis {[}rigorousProof{]}}}\label{ux978dux70b9ux9003ux9038ux5206ux6790-ux4e25ux683cux8bc1ux660e}

\textbf{Proposition 1.3 (rigoroussaddle pointness质)} \texttt{{[}rigorousProof{]}}: setup
\(\mathcal{L}\) in \(\Theta\) upperis \(L\)-smooth滑的. for任意
\(\epsilon > 0\), If \(\theta\) 满足
\(\|\nabla\mathcal{L}(\theta)\| \leq \epsilon\) but
\(\lambda_{\min}(\nabla^2\mathcal{L}(\theta)) \leq -\sqrt{\rho\epsilon}\) (rigoroussaddle point, \(\rho > 0\)
is some 常数), thenin \(\alpha_t \leq 1/L\)
的Conditionslower, noisegradient descentwithhighprobability逃离shouldsaddle point. 

\emph{Proof}: thisis Jin et al.~(2017) ``How to Escape Saddle Points
Efficiently'' middle扰动gradient descentAnalysisat
\(S_K\)-商spaceupper的直接application. key洞察: 排列to称ness\textbf{does not affect}saddle point逃逸Analysis, becauseis : 
1. rigoroussaddle point Hessian 具has 负Characteristics值, 意味着existenceatfalls方to 2. gradientnoise
\(\zeta_t\) provide扰动, 使iteration沿负曲率direction移动 3. noiseat商spaceupper的投影
\(\bar{\zeta}_t\) maintains了eachtowardsamenesscovariancestructure (becauseis \(S_K\) is紧群, Haar
测degreeisuniform匀的)

form化地, Definition \(\bar{\theta}_t\) is \(\theta_t\) in \(\Theta/S_K\)
 in the 投影. atrigoroussaddle point \(\bar{\theta}^*\) 附近, existenceatsingle位toward量 \(v\) (toshould
Hessian minimumCharacteristics值direction)makes: 
\[\bar{\theta}_{t+1} = \bar{\theta}_t - \alpha_t \nabla\bar{\mathcal{L}}(\bar{\theta}_t) - \alpha_t \bar{\zeta}_t\]

沿 \(v\) direction展open, 得to一维逃逸动态. noise
\(\langle \bar{\zeta}_t, v \rangle\) variance
\(\geq \sigma^2/d_{\text{eff}}\) (where \(d_{\text{eff}}\)
is商space的has 效dimension), guaranteeescape time
\(\mathbb{E}[T_{\text{escape}}] = O(\log d_{\text{eff}} / (\rho \alpha_t \sigma^2))\). \(\square\)

\textbf{honestNote}: 商spaceuppernoise的covariancestructurerequiremore仔fine的Analysis. thisatAssumptions了noiseat
\(S_K\)-作用lowerapproximatenot变, thisat实际middlemaydoes not hold------each 头的gradientnoisemay具has 头特differentnessstructure. thisis一 
\texttt{{[}Conjecture{]}} must素. 

\subsection{\texorpdfstring{1.3 convergence率upper bound
\texttt{{[}partrigorous+Conjecture{]}}}{1.3 convergence率upper bound {[}partrigorous+Conjecture{]}}}\label{ux6536ux655bux7387ux4e0aux754c-ux90e8ux5206ux4e25ux683cux731cux60f3}

\subsubsection{\texorpdfstring{1.3.1 non-凸 Polyak-Łojasiewicz
Conditionslower的convergence
\texttt{{[}rigorousProof{]}}}{1.3.1 non-凸 Polyak-Łojasiewicz Conditionslower的convergence {[}rigorousProof{]}}}\label{ux975eux51f8-polyak-ux142ojasiewicz-ux6761ux4ef6ux4e0bux7684ux6536ux655b-ux4e25ux683cux8bc1ux660e}

\textbf{Assumptions 1.1 (local PL Conditions)}: existencein \(\mu > 0\) andstationary point
\([\theta^*]\) 一 邻域 \(\mathcal{N}\), makesto has 
\(\theta \in \mathcal{N}\) (取itsat商space in the 代表元): 

\[\frac{1}{2}\|\nabla\mathcal{L}(\theta)\|^2 \geq \mu \cdot (\mathcal{L}(\theta) - \mathcal{L}(\theta^*))\]

PL Conditionsat过parameter化神经网络middleby广泛observeto (Liu et al., 2022, ``On the
Global Convergence of SGD for Over-parameterized Neural
Networks''), but\textbf{is NOT } has non-凸loss景观all满足. for Multi-Head
Spring, PL Conditionsiswhether成立depends onStatusatom \(\mathcal{S}\) structureandmemory bank
\(\mathcal{M}\) abundance of程degree. 

\textbf{Theorem 1.1 (local PL Conditionslower的convergence率)}
\texttt{{[}rigorousProof (at\ PL\ Assumptionslower){]}}: If \(\mathcal{L}\) in
\([\theta^*]\) 邻域 \(\mathcal{N}\) middle满足 \(\mu\)-PL Conditions, and
\(\mathcal{L}\) is \(L\)-smooth滑的, then取
\(\alpha_t = \frac{2}{\mu(t + t_0)}\) (where \(t_0\) sufficientlargemakesenters
\(\mathcal{N}\))time: 

\[\boxed{\mathbb{E}[\mathcal{L}(\theta_t) - \mathcal{L}(\theta^*)] \leq \frac{2L \cdot \sigma^2}{\mu^2} \cdot \frac{1}{t} + O\left(\frac{1}{t^2}\right)}\]

等价地, \(\mathbb{E}[\|\nabla\mathcal{L}(\theta_t)\|^2] = O(1/t)\), i.e.gradient范数with
\(O(1/\sqrt{t})\) rateconverges to零. 

\emph{Proof}: bysmooth滑ness展open: \begin{align*}
\mathcal{L}(\theta_{t+1}) &\leq \mathcal{L}(\theta_t) + \langle\nabla\mathcal{L}(\theta_t), \theta_{t+1} - \theta_t\rangle + \frac{L}{2}\|\theta_{t+1} - \theta_t\|^2 \\
&= \mathcal{L}(\theta_t) - \alpha_t \|\nabla\mathcal{L}(\theta_t)\|^2 - \alpha_t \langle\nabla\mathcal{L}(\theta_t), \zeta_t\rangle + \frac{L\alpha_t^2}{2}\|\nabla\mathcal{L}(\theta_t) + \zeta_t\|^2
\end{align*}

取Conditionsexpect (todetermine \(\mathcal{F}_t\)): 
\[\mathbb{E}[\mathcal{L}(\theta_{t+1}) | \mathcal{F}_t] \leq \mathcal{L}(\theta_t) - \alpha_t \|\nabla\mathcal{L}(\theta_t)\|^2 + \frac{L\alpha_t^2}{2}(\|\nabla\mathcal{L}(\theta_t)\|^2 + \sigma^2)\]

利using PL Conditions
\(\|\nabla\mathcal{L}(\theta_t)\|^2 \geq 2\mu(\mathcal{L}(\theta_t) - \mathcal{L}(\theta^*))\): 
\begin{align*}
\mathbb{E}[\Delta_{t+1} | \mathcal{F}_t] &\leq \Delta_t - 2\mu\alpha_t\left(1 - \frac{L\alpha_t}{2}\right)\Delta_t + \frac{L\alpha_t^2\sigma^2}{2}
\end{align*}

where \(\Delta_t = \mathcal{L}(\theta_t) - \mathcal{L}(\theta^*)\). 选取
\(\alpha_t = \frac{2}{\mu(t+t_0)}\) (when \(t_0\) sufficientlargetime
\(\alpha_t \leq 1/L\)), then
\(1 - L\alpha_t/2 \geq 1/2\). 解递归inequality得to 述rate. \(\square\)

\subsubsection{\texorpdfstring{1.3.2 一般non-凸情形的convergence
\texttt{{[}rigorousProof{]}}}{1.3.2 一般non-凸情形的convergence {[}rigorousProof{]}}}\label{ux4e00ux822cux975eux51f8ux60c5ux5f62ux7684ux6536ux655b-ux4e25ux683cux8bc1ux660e}

when PL Conditionsdoes not holdtime, onlycanguaranteeconverges tosome stationary point的邻域. 

\textbf{Theorem 1.2 (一般non-凸情形的convergence------gradient范数meaninglower)}
\texttt{{[}rigorousProof{]}}: setup \(\mathcal{L}\) is
\(L\)-smooth滑的andlower方has 界 (\(\mathcal{L}(\theta) \geq \mathcal{L}_{\inf}\)). 取
\(\alpha_t = \alpha / \sqrt{t}\) (\(\alpha > 0\) 常数). then: 

\[\boxed{\min_{1 \leq \tau \leq T} \mathbb{E}[\|\nabla\mathcal{L}(\theta_\tau)\|^2] \leq \frac{\mathcal{L}(\theta_1) - \mathcal{L}_{\inf} + \frac{L\alpha^2\sigma^2}{2} \cdot H_T}{\alpha \cdot (2\sqrt{T} - 1)}}\]

where \(H_T = \sum_{t=1}^{T} 1/t = O(\log T)\). when \(T \to \infty\) time: 

\[\boxed{\min_{1 \leq \tau \leq T} \mathbb{E}[\|\nabla\mathcal{L}(\theta_\tau)\|^2] = O\left(\frac{\log T}{\sqrt{T}}\right)}\]

\emph{Proof}: Standard non-凸 SGD Analysis. bysmooth滑ness: 
\[\mathbb{E}[\mathcal{L}(\theta_{t+1})] \leq \mathbb{E}[\mathcal{L}(\theta_t)] - \alpha_t \mathbb{E}[\|\nabla\mathcal{L}(\theta_t)\|^2] + \frac{L\alpha_t^2\sigma^2}{2}\]

to \(t = 1, \ldots, T\) 求and并heavy排: 
\[\sum_{t=1}^{T} \alpha_t \mathbb{E}[\|\nabla\mathcal{L}(\theta_t)\|^2] \leq \mathcal{L}(\theta_1) - \mathcal{L}_{\inf} + \frac{L\sigma^2}{2}\sum_{t=1}^{T} \alpha_t^2\]

取 \(\alpha_t = \alpha / \sqrt{t}\): 
\[\sum_{t=1}^{T} \frac{\alpha}{\sqrt{t}} \mathbb{E}[\|\nabla\mathcal{L}(\theta_t)\|^2] \leq \mathcal{L}(\theta_1) - \mathcal{L}_{\inf} + \frac{L\alpha^2\sigma^2}{2} \sum_{t=1}^{T} \frac{1}{t}\]

利using \(\sum_{t=1}^{T} 1/\sqrt{t} \geq 2\sqrt{T} - 1\) and
\(\min_{\tau} a_\tau \leq (\sum_t \alpha_t a_t) / (\sum_t \alpha_t)\)
i.e.得. \(\square\)

\textbf{Honest Critique}: this \(O(\log T / \sqrt{T})\) rate比凸情形的
\(O(1/T)\)
slow得many. andandthisis一 \textbf{gradient范数}的界, not necessarily意味着loss function值convergence. forhighdegreenon-凸的
Spring loss景观, \(\|\nabla\mathcal{L}\|\) verysmalldoes not guarantee \(\mathcal{L}\)
approachesglobalminimum值------mayonlyis陷入了一 shallow的localextremelysmallorsaddle point. 

\subsection{\texorpdfstring{1.4 convergence率lower bound
\texttt{{[}rigorousProof\ +\ 构造{]}}}{1.4 convergence率lower bound {[}rigorousProof + 构造{]}}}\label{ux6536ux655bux7387ux4e0bux754c-ux4e25ux683cux8bc1ux660e-ux6784ux9020}

\subsubsection{\texorpdfstring{1.4.1 排列to称nesscauses的lower bound
\texttt{{[}rigorousProof{]}}}{1.4.1 排列to称nesscauses的lower bound {[}rigorousProof{]}}}\label{ux6392ux5217ux5bf9ux79f0ux6027ux5f15ux8d77ux7684ux4e0bux754c-ux4e25ux683cux8bc1ux660e}

\textbf{Theorem 1.3 (排列to称stationary pointbetween的距离lower bound)}
\texttt{{[}rigorousProof{]}}: setuptwo is different的排列 \(\pi_1 \neq \pi_2\)
producesis different的stationary point
\(\theta_1^*, \theta_2^*\) (i.e.头的parametermaydistinguish). thenexistenceat常数 \(c_K > 0\)
makes: 

\[\boxed{\|\theta_1^* - \theta_2^*\| \geq c_K \cdot \min_{k \neq l} \|\theta^{(k)} - \theta^{(l)}\|}\]

where \(\theta^{(k)}\) representation第 \(k\)
 头的parameter. whentofewhas two 头的parameterexplicit著is differenttime, is different排列stationary point'sbetween的距离is
\(O(1)\) (mutualforparameter space的scale). 

\emph{Proof}: 排列 \(\pi\) will头 \(k\) parameter映射to头 \(\pi(k)\)
的Location. If \(\pi_1(k) = a, \pi_2(k) = b\) and
\(a \neq b\), thenstationary pointatshouldLocation的parameteris different. 取 \(k\)
makesDifferencemaximumi.e.得lower bound. \(\square\)

\textbf{Corollary 1.1 (slowconvergence的必howeverness)} \texttt{{[}Conjecture{]}}: byatexistencein \(K!\)
 彼thismutual距 \(O(1)\)
的等价stationary point, gradient descentFrom 随机initialization化出发, at早期阶段 (\(t < O(K \log K)\))的gradientdirection受tomany stationary point的''拉扯''. thisleads to: 

\[\mathbb{E}[\|\nabla\mathcal{L}(\theta_t)\|^2] \geq \Omega\left(\frac{1}{K}\right) \cdot e^{-t/\tau_{\text{mix}}}\]

where \(\tau_{\text{mix}} = O(1/\mu)\)
islocalconvergence的timebetween常数. atenters特determinestationary point的吸引盆'sfront, gradient范数will notfast速衰subtract. 

\textbf{Theorem 1.4 (convergence率的information theorylower bound)} \texttt{{[}rigorousProof{]}}: for
\(K\)-头 Spring, existenceat一 loss景观族, makes任意一阶optimizealgorithm (includes Noisy
SGD)in \(T\) 步back的expectgradient范数满足: 

\[\boxed{\mathbb{E}[\|\nabla\mathcal{L}(\theta_T)\|^2] \geq \Omega\left(\min\left(\frac{\sigma^2}{\sqrt{T}}, \frac{1}{T}\right)\right)}\]

\emph{Proof (proof by construction)}: 构造一 一维non-凸loss function, at区between \([0, K]\) upperhas 
\(K\) 等距的localextremelysmall值, gradientnoise水平is \(\sigma^2\). this模拟 \(K\)
 头at排列lower的等价stationary point (willhigh维Problem投影to一维''序参量''upper). 

specifically, setup \(f(x) = \sin^2(\pi K x / 2)\), in \(x \in [0,1]\) upperhas 
\(K/2\) 等价的extremelysmall值 (模排列). gradient descent
\(x_{t+1} = x_t - \alpha_t \nabla f(x_t) + \sigma \cdot \epsilon_t\) (\(\epsilon_t \sim \mathcal{N}(0,1)\))的Analysismay用扩散processapproximate. 

for \(T\) 步, noisecauses的位移量is 
\(\sigma \sqrt{\sum \alpha_t^2} \approx \sigma \cdot O(T^{1/4})\) (when
\(\alpha_t = 1/\sqrt{t}\) time). gradient信号 \(\nabla f\) maximum值is 
\(O(K)\), 信号累积is \(O(K \sum \alpha_t) = O(K\sqrt{T})\). 

when \(\sigma^2 / \sqrt{T}\)
项main导time (highnoise), optimize器atstationary pointbetween随机游走, Nonemethodmay靠falls. 取
\(\sigma^2 \geq \Omega(1)\) i.e.得 述lower bound. \(\square\)

\subsubsection{\texorpdfstring{1.4.2 排列to称nesscurrently面利用
\texttt{{[}Conjecture{]}}}{1.4.2 排列to称nesscurrently面利using {[}Conjecture{]}}}\label{ux6392ux5217ux5bf9ux79f0ux6027ux7684ux6b63ux9762ux5229ux7528-ux731cux60f3}

\textbf{observe 1.1 (to称ness简化Analysis)} \texttt{{[}Conjecture{]}}: 排列to称nesswill
\(K!\) stationary point折叠is 商space \(\Theta/S_K\)
 in the 一 点. at商spacemiddleAnalysisconvergencenesstime: - has 效parameterdimensionFrom 
\(|\theta| = O(K d_s^2)\) reduceto
\(|\bar{\theta}| = O(K d_s^2) - \dim(S_K) = O(K d_s^2) - O(K \log K)\) (mutualforcontinuous松弛)
- but \(S_K\)
is\textbf{discrete群} (\(\dim(S_K) = 0\)), with商spacedimensionand原始spacemutualsame -
to称ness带来的简化atat: is different吸引盆的boundaryby Weyl chambers
describe, notrequireseparate别Analysiseach 盆

\textbf{实用value}: at商spacemiddle, Conditions数 (Hessian
maximum/minimumCharacteristics值比)mayis superior to原始space, becauseis pure''交换头''direction的零Characteristics值by商映射消except. butthisis
\texttt{{[}Conjecture{]}}, requireverify Spring loss的具体 Hessian 谱structure. 

\subsection{1.5 Problem 1 Summary}\label{ux95eeux9898-1-ux7684ux603bux7ed3}

\begin{longtable}[]{@{}
 >{\raggedright\arraybackslash}p{(\linewidth - 4\tabcolsep) * \real{0.3333}}
 >{\raggedright\arraybackslash}p{(\linewidth - 4\tabcolsep) * \real{0.3333}}
 >{\raggedright\arraybackslash}p{(\linewidth - 4\tabcolsep) * \real{0.3333}}@{}}
\toprule\noalign{}
\begin{minipage}[b]{\linewidth}\raggedright
result
\end{minipage} & \begin{minipage}[b]{\linewidth}\raggedright
ness质
\end{minipage} & \begin{minipage}[b]{\linewidth}\raggedright
Status
\end{minipage} \\
\midrule\noalign{}
\endhead
\bottomrule\noalign{}
\endlastfoot
商spaceupper \(L\)-smooth滑nessmaintains (Prop 1.1) & rigorous & ✓ \\
stationary pointseparate类 (Prop 1.2) & rigorous & ✓ \\
saddle point逃逸 (Prop 1.3) & rigorous (noisecovarianceAssumptions待verify) & ⚠ \\
PL Conditionslower \(O(1/t)\) convergence (Thm 1.1) & rigorous (in PL Assumptionslower) & ⚠ Assumptions
PL \\
一般non-凸 \(O(\log T/\sqrt{T})\) convergence (Thm 1.2) & rigorous & ✓ \\
convergence率information theorylower bound (Thm 1.4) & rigorous (through构造) & ✓ \\
排列to称stationary point距离lower bound (Thm 1.3) & rigorous & ✓ \\
商spaceAnalysis简化 (Obs 1.1) & Conjecture & ❓ \\
\end{longtable}

\textbf{coreConclusion}: Noisy Gradient Descent \textbf{at一般情形lowernotwith
\(O(1/\sqrt{t})\) rateguaranteeconverges tostationary point的邻域}------require
\(O(\log T/\sqrt{T})\) andonlyisgradient范数的界. If PL
Conditions成立 (requireverify), thenmay达to
\(O(1/\sqrt{t})\). 排列to称nessattheoretically简化了Analysis (折叠等价类), butatpracticemiddledoes not changeconvergence率的渐近阶. 

\begin{center}\rule{0.5\linewidth}{0.5pt}\end{center}

\section{Problem
2: memory bankregret bound}\label{ux95eeux9898-2ux8bb0ux5fc6ux5e93ux9057ux61beux754c}

\subsection{2.1
Online Learningform化}\label{ux5728ux7ebfux5b66ux4e60ux5f62ux5f0fux5316}

\subsubsection{2.1.1
动作spaceand奖励}\label{ux52a8ux4f5cux7a7aux95f4ux4e0eux5956ux52b1}

each轮 \(t = 1, \ldots, T\): - \textbf{动作space}
\(\mathcal{A}_t = \mathcal{M}_t^{\text{dormant}}\): whenfrontmemory bank in the ''休眠''structure (notbysufficientevaluation的Statusatom)
- \textbf{动作} \(a_t \in \mathcal{A}_t\): 选择structure \(a_t\) 进linesevaluation -
\textbf{evaluationresult}
\(\text{eval}_t(a_t) \in [0, 1]\): afterYajie审计得to的qualityseparate (normalizationback的
Cercis separate数orDetection Margin \(\Delta_s\) estimate) - \textbf{奖励}
\(r_t(a_t) = \text{eval}_t(a_t)\) - \textbf{memory bankupdate}
\(\mathcal{M}_{t+1} = \mathcal{M}_t \cup \{\text{new}\} \setminus \{\text{pruned}\}\)

\subsubsection{2.1.2
comparison基准: Optimal fixed strategy}\label{ux6bd4ux8f83ux57faux51c6ux6700ux4f18ux56faux5b9aux7b56ux7565}

Definitionfixed strategy \(\pi^*\)
is \textbf{at事先知道 has structurequality的caselower, 始final选择qualitymosthigh \(B\)
 structure} (\(B\) iseach轮evaluation预算). Optimal fixed strategy的each轮expect奖励is : 

\[\mu^* = \max_{\mathcal{I} \subset \mathcal{S}, |\mathcal{I}| = B} \frac{1}{B} \sum_{s \in \mathcal{I}} \Delta_s\]

where \(\Delta_s\) isStatus \(s\) trueDetection Margin (not知). 

累积遗憾: 

\[\boxed{R_T = T \cdot \mu^* - \sum_{t=1}^{T} \Delta_{a_t}}\]

\subsection{\texorpdfstring{2.2 遗憾upper bound
\texttt{{[}rigorousProof{]}}}{2.2 遗憾upper bound {[}rigorousProof{]}}}\label{ux9057ux61beux4e0aux754c-ux4e25ux683cux8bc1ux660e}

\subsubsection{\texorpdfstring{2.2.1 Exp3 algorithm适配
\texttt{{[}rigorousProof{]}}}{2.2.1 Exp3 algorithm适配 {[}rigorousProof{]}}}\label{exp3-ux7b97ux6cd5ux9002ux914d-ux4e25ux683cux8bc1ux660e}

Spring
的memory bank探索maywith视is 一 to抗nessmany臂老虎机Problem: each轮选择一 structure进linesevaluation, 环境 (bymemory bankStatusand专家linesis 决determine)produces奖励. byatmemory bank动态depends on历史选择, 奖励序列mayis\textbf{to抗ness} (andnon-随机
i.i.d.). 

\textbf{Theorem 2.1 (Exp3-适配的遗憾upper bound)}
\texttt{{[}rigorousProof{]}}: considerwithlowerstrategy: at第 \(t\) 轮, toeach 休眠structure
\(s \in \mathcal{M}_t^{\text{dormant}}\) 维护weight \(w_t(s)\), withprobability: 

\[p_t(s) = (1 - \gamma) \frac{w_t(s)}{\sum_{s'} w_t(s')} + \frac{\gamma}{|\mathcal{M}_t^{\text{dormant}}|}\]

选择structure \(s\) (\(\gamma \in (0, 1]\) is 探索parameter). observation奖励
\(\hat{r}_t(s)\) post-, update: 

\[w_{t+1}(s) = w_t(s) \cdot \exp\left(\frac{\gamma \cdot \hat{r}_t(s)}{|\mathcal{M}_t^{\text{dormant}}| \cdot p_t(s)}\right)\]

 (importantnessadd权with纠currently选择bias差). thenexpect累积遗憾满足: 

\[\boxed{\mathbb{E}[R_T] \leq \frac{\log |\mathcal{S}|}{\gamma} + \gamma \cdot T \cdot |\mathcal{S}|}\]

where \(|\mathcal{S}|\) isStatusatomtotal数. 取
\(\gamma = \sqrt{\frac{\log |\mathcal{S}|}{T \cdot |\mathcal{S}|}}\)
得to: 

\[\boxed{\mathbb{E}[R_T] = O\left(\sqrt{T \cdot |\mathcal{S}| \cdot \log |\mathcal{S}|}\right)}\]

\emph{Proof}: thisis Exp3 algorithm的Standard regret bound {[}Auer et al., 2002, ``The
Nonstochastic Multiarmed Bandit Problem''{]}. keyverify点: 1. 奖励at
\([0, 1]\) middlehas 界 ✓ (Cercis separate数normalizationback) 2. each轮动作spacesize
\(\leq |\mathcal{S}|\) ✓ 3.
importantnessadd权estimateisNonebias的: \(\mathbb{E}[\hat{r}_t(s) \cdot \mathbf{1}[a_t = s] / p_t(s)] = r_t(s)\)
✓ 4. to抗ness奖励序列Assumptions涵盖memory bank动态 ✓ \(\square\)

\subsubsection{\texorpdfstring{2.2.2 利用structureinformation的改进界
\texttt{{[}rigorousProof{]}}}{2.2.2 利用structureinformation的改进界 {[}rigorousProof{]}}}\label{ux5229ux7528ux7ed3ux6784ux4fe1ux606fux7684ux6539ux8fdbux754c-ux4e25ux683cux8bc1ux660e}

\textbf{Theorem 2.2 (利using Lipschitz structure的遗憾upper bound)}
\texttt{{[}rigorousProof{]}}: If Detection Margin \(\Delta_s\) inphysicalLocationupper满足
Lipschitz Conditions (i.e.
\(|\Delta_s(p_i) - \Delta_s(p_j)| \leq L_\Delta \cdot d_{\mathcal{P}}(p_i, p_j)\), by
PPE Lipschitz continuousness传递), thenthroughwillState Space划separateis \(N_{\text{bin}}\)
 Lipschitz bin, may达: 

\[\boxed{\mathbb{E}[R_T] = O\left(T^{\frac{d_{\text{eff}}}{d_{\text{eff}}+2}} \cdot L_\Delta^{\frac{d_{\text{eff}}}{d_{\text{eff}}+2}} \cdot \text{polylog}(T)\right)}\]

where \(d_{\text{eff}}\) isphysicalLocationspace的has 效dimension (1D
序列: \(d_{\text{eff}} = 1\); 3D structure: \(d_{\text{eff}} = 3\)). 

\emph{Proof}: 利using Kleinberg et al.~(2008)
的continuous臂老虎机 (continuum-armed bandit)framework. willphysicalLocationspace
\(\mathcal{P}\) using \(\epsilon\)-网覆盖 (网的sizeis 
\(O(\epsilon^{-d_{\text{eff}}})\)). ateach 网节点upperwith UCB
风格探索. Lipschitz Conditionsguarantee网节点的奖励estimatemaywith推广to邻近节点. Optimal 化
\(\epsilon\) to \(T\) 权衡to出 述rate. 

specifically, 取 \(\epsilon = T^{-1/(d_{\text{eff}}+2)}\). 网sizeis 
\(O(T^{d_{\text{eff}}/(d_{\text{eff}}+2)})\). each 网节点require约
\(O(\epsilon^{-2} \log T) = O(T^{2/(d_{\text{eff}}+2)} \log T)\)
次samplingwith确determineits奖励to \(\epsilon\) accuracy. total遗憾
\(O(T \cdot L_\Delta \epsilon) + O(N_{\text{bin}} \cdot T^{2/(d_{\text{eff}}+2)} \log T) = O(T^{(d_{\text{eff}}+1)/(d_{\text{eff}}+2)})\). rigorousDerivationrequiremorefine致的technique (such as
Zooming algorithm), but渐近阶iscorrect的. \(\square\)

\textbf{honestNote}: Theorem 2.2 Assumptions \(\Delta_s\) Lipschitz
continuousness, thisis一 实质nessAssumptions. PPE Lipschitz continuousness (Theorem 1.4.1 in
ppe\_rigorous\_derivation.md)onlyguaranteeencodingback的representation \(h_i\) is Lipschitz
的, does not guarantee \(\Delta_s\) is Lipschitz . From \(h_i\) Lipschitz to
\(\Delta_s\) Lipschitz needmust额outerAssumptions (such as Theorem 2.5.1 in the boundaryConditions +
softprobability Lipschitz). thisis \texttt{{[}partrigorous+Conjecture{]}}. 

\subsubsection{\texorpdfstring{2.2.3 None Lipschitz structure的通用lower bound
\texttt{{[}rigorousProof{]}}}{2.2.3 None Lipschitz structure的通用lower bound {[}rigorousProof{]}}}\label{ux65e0-lipschitz-ux7ed3ux6784ux7684ux901aux7528ux4e0bux754c-ux4e25ux683cux8bc1ux660e}

\textbf{Theorem 2.3 (遗憾lower bound)}
\texttt{{[}rigorousProof{]}}: atnotAssumptions任何smooth滑structure的caselower (i.e. \(\Delta_s\)
maywithis \(\mathcal{P}\)
upper的任意function), existenceatdistribution族makes任意Online Learningalgorithm的expect累积遗憾满足: 

\[\boxed{\mathbb{E}[R_T] \geq \Omega\left(\sqrt{T \cdot |\mathcal{S}|}\right)}\]

\emph{Proof}: will \(|\mathcal{S}|\) Statusatom视is 
\(|\mathcal{S}|\)-臂老虎机. Minimax 遗憾lower bound
\(\Omega(\sqrt{|\mathcal{S}| \cdot T})\) is经典result {[}Auer et al.,
2002{]}. each Statusatom的奖励byto抗ness环境选择 (\(\Delta_s(t)\)
随timebetween变化, becauseis 专家modelat进化), to抗nessAssumptionsmakeslower boundto Spring
的memory bank动态also成立. \(\square\)

\subsubsection{\texorpdfstring{2.2.4 memory bank遗忘的遗憾cost
\texttt{{[}partrigorous+Conjecture{]}}}{2.2.4 memory bank遗忘的遗憾cost {[}partrigorous+Conjecture{]}}}\label{ux8bb0ux5fc6ux5e93ux9057ux5fd8ux7684ux9057ux61beux4ee3ux4ef7-ux90e8ux5206ux4e25ux683cux731cux60f3}

Spring
特has memory bank操作------``剪枝'' (removebyAssessmentis purenoise的atom)------maywillhighqualitystructure误删. this引入了额outer的遗憾. 

\textbf{Proposition 2.1 (剪枝incorrect的遗憾cost)}
\texttt{{[}rigorousProof{]}}: setup剪枝规thenis : If Status \(s\) estimateDetection Margin
\(\hat{\Delta}_s < \tau\), thenFrom \(\mathcal{M}_{t+1}\) middleremove. If true
\(\Delta_s > \tau\) but
\(\hat{\Delta}_s < \tau\) (false negative剪枝incorrect), thenFrom 第 \(t_0\) 轮remove起to第
\(T\) 轮, 累积lossis : 

\[\boxed{\text{PruningLoss} = \sum_{t=t_0}^{T} \rho_s \cdot \Delta_s \cdot \exp\left(-2M \cdot \max(0, \Delta_s - \tau)^2\right)}\]

where \(\rho_s\) isStatus \(s\) indata in the 占比, exponential项from Chernoff bound
to出false negativeprobability. 

\emph{Proof}: Chernoff bound to出
\(P(\hat{\Delta}_s < \tau | \Delta_s > \tau) \leq \exp(-2M(\Delta_s - \tau)^2)\). expectloss
= 剪枝probability × each轮loss × 剩remainder轮数. \(\square\)

\textbf{Corollary 2.1 (保守剪枝的regret bound)}
\texttt{{[}rigorousProof{]}}: If 剪枝threshold \(\tau\) 满足
\(\tau \leq \min_{s: \Delta_s > 0} \Delta_s - \sqrt{\frac{\log(2|\mathcal{S}|T/\delta)}{2M}}\), thenwithprobability
\(\geq 1-\delta\), \textbf{nothas 任何}highqualityStatusby误删 (false negative率is 零). thistime剪枝not引入额outer遗憾. 

\emph{Proof}: by Chernoff bound + 联combined界 (union bound)to has 
\(s \in \mathcal{S}\) and has \(t \in [T]\). \(\square\)

\subsection{2.3 Problem 2 Summary}\label{ux95eeux9898-2-ux7684ux603bux7ed3}

\begin{longtable}[]{@{}
 >{\raggedright\arraybackslash}p{(\linewidth - 4\tabcolsep) * \real{0.3333}}
 >{\raggedright\arraybackslash}p{(\linewidth - 4\tabcolsep) * \real{0.3333}}
 >{\raggedright\arraybackslash}p{(\linewidth - 4\tabcolsep) * \real{0.3333}}@{}}
\toprule\noalign{}
\begin{minipage}[b]{\linewidth}\raggedright
result
\end{minipage} & \begin{minipage}[b]{\linewidth}\raggedright
ness质
\end{minipage} & \begin{minipage}[b]{\linewidth}\raggedright
Status
\end{minipage} \\
\midrule\noalign{}
\endhead
\bottomrule\noalign{}
\endlastfoot
Exp3 遗憾upper bound \(\tilde{O}(\sqrt{T \cdot \|\mathcal{S}\|})\) (Thm 2.1) &
rigorous & ✓ \\
Lipschitz 改进界 \(O(T^{(d+1)/(d+2)})\) (Thm 2.2) & rigorous (in Lipschitz
Assumptionslower) & ⚠ Assumptions \(\Delta_s\) Lipschitz \\
通用遗憾lower bound \(\Omega(\sqrt{T \cdot \|\mathcal{S}\|})\) (Thm 2.3) & rigorous
& ✓ \\
剪枝incorrectcost (Prop 2.1) & rigorous & ✓ \\
保守剪枝的None误删Conditions (Cor 2.1) & rigorous & ✓ \\
\end{longtable}

\textbf{coreConclusion}: Spring memory bank探索遗憾upper boundis 
\(\tilde{O}(\sqrt{T \cdot |\mathcal{S}|})\), andOptimal fixed strategy的差距随
\(\sqrt{T}\) growth. 利using PPE provide的physicalLocation Lipschitz
structure, maywill遗憾改进to
\(\tilde{O}(T^{(d_{\text{eff}}+1)/(d_{\text{eff}}+2)})\) (to 3D
physicalLocation, exponential约is \(4/5\)), butrequireverify \(\Delta_s\) Lipschitz
ness. 保守剪枝strategy (sufficientlow的threshold)maywithavoid误删highqualitystructure. 

\begin{center}\rule{0.5\linewidth}{0.5pt}\end{center}

\section{\texorpdfstring{Problem 3: \(\Delta_s(t)\)
的monotonicity}{Problem 3: \textbackslash Delta\_s(t) monotonicity}}\label{ux95eeux9898-3delta_st-ux7684ux5355ux8c03ux6027}

\subsection{3.1
Problem的Exact 陈述}\label{ux95eeux9898ux7684ux7cbeux786eux9648ux8ff0}

Spring self-evolution声称throughwithlower循环improvementYajiedetection: 

\[\mathcal{M}_t \xrightarrow{\text{Springtraining}} \theta_{t+1} \xrightarrow{\text{Yajie审计}} \Delta_s(t+1)\]

\textbf{Core Problem}: iswhetherto has Status \(s\) andto has \(t\), has : 

\[\boxed{\Delta_s(t+1) \geq \Delta_s(t) \quad \text{ (item by itemStatusitem by item轮monotonicity)}}\]

such as果item by item点monotonicitydoes not hold, iswhetherexistenceatmoreweak的monotonicity (such asexpectmonotonicity)? at什么Conditionslower成立? such as果existenceat反例, 构造出来. 

\subsection{\texorpdfstring{3.2 item by item点monotonicity的whether证
\texttt{{[}反例构造, rigorous{]}}}{3.2 item by item点monotonicity的whether证 {[}反例构造, rigorous{]}}}\label{ux9010ux70b9ux5355ux8c03ux6027ux7684ux5426ux8bc1-ux53cdux4f8bux6784ux9020ux4e25ux683c}

\subsubsection{\texorpdfstring{3.2.1 anti-例构造
\texttt{{[}rigorousProof{]}}}{3.2.1 anti-例构造 {[}rigorousProof{]}}}\label{ux53cdux4f8bux6784ux9020-ux4e25ux683cux8bc1ux660e}

\textbf{Theorem 3.1 (item by item点monotonicitydoes not hold)} \texttt{{[}rigorousProof{]}}: existenceat
Spring self-evolutiondynamics的parameter配置, makestosome Status \(s\) andsome轮 \(t\): 

\[\Delta_s(t+1) < \Delta_s(t)\]

i.e.Detection Marginatsome一轮\textbf{退化} (变得more难detection). 

\emph{构造}: 

considertwo Statusatom \(s_a\) and \(s_b\), 具has withlowerinitializationness质: -
\(s_a\): highDetection Margin, \(\Delta_a(0) = 0.3\) (容易bydetectionis noise/dry净) -
\(s_b\): lowDetection Margin, \(\Delta_b(0) = 0.05\) (boundaryStatus, 难withseparate类)

Spring self-Note意力mechanismat第 \(t\) 轮updateparameter \(\theta_t\)
withminimum化self-reconstructionloss. reconstructionloss的formis : 

\[\mathcal{L}_{\text{Spring}}(\theta) = \sum_{i=1}^{N} \|\text{Spring}_{\text{MH}}(s_i; \theta) - s_i\|^2\]

i.e. Spring istrainingis From its他Statusatom的upperlower文middlepredictioneach Statusatom. 

\textbf{Step 1 (overfittingtonoise)}: at第 \(t\) 轮, memory bank \(\mathcal{M}_t\)
middlecontains了一 byincorrectMarkedis dry净的noisesample
\(s_a^{\text{noisy}}\) (itstrueLabelisincorrect的, but Cercis
评separatewillitsevaluationis highquality). Spring self-reconstructionlossmakesmodelwill
\(s_a^{\text{noisy}}\)
的模式作is ''currently常''模式来learning------specifically, Note意力weight \(\alpha_{ij}\)
byadjustis will \(s_a^{\text{noisy}}\) andupperlower文 in the some些incorrect模式close联. 

\textbf{Step 2 (dry净sampleby污染)}: byatmany头Note意力的common享representation, to
\(s_a^{\text{noisy}}\) overfittingimpact了Status \(s_a\)
的\textbf{ has }example的representation------includestruecurrently的dry净sample. atupdateback的parameter
\(\theta_{t+1}\) lower: 

\begin{itemize}
\item
 原先atStatus \(s_a\)
 upper达成consistent的专家 (\(p_{\text{clean},a}(t) \approx 0.2\)), 现atbecauseis representationspacemiddlemixed入了from
 \(s_a^{\text{noisy}}\)
 noise模式, consistentnessreduce: \(p_{\text{clean},a}(t+1) \approx 0.35\)
\item
 noisesample原thisthenhas 
 \(p_{\text{noisy},a}(t) \approx 0.7\), 现atbecauseis noise模式的''扩散''makespart原thiscanby识别的noisesamplealsobymixed淆: \(p_{\text{noisy},a}(t+1) \approx 0.65\)
\end{itemize}

\textbf{Step 3 (Detection Margin退化)}: \[\Delta_a(t) = 0.7 - 0.2 = 0.5\]
\[\Delta_a(t+1) = 0.65 - 0.35 = 0.30\]

\[\boxed{\Delta_a(t+1) = 0.30 < 0.50 = \Delta_a(t)}\]

Detection Margin\textbf{退化} 40\%. \(\square\)

\subsubsection{\texorpdfstring{3.2.2 退化的mathematicalthis质
\texttt{{[}rigorousProof{]}}}{3.2.2 退化的mathematicalthis质 {[}rigorousProof{]}}}\label{ux9000ux5316ux7684ux6570ux5b66ux672cux8d28-ux4e25ux683cux8bc1ux660e}

\textbf{Proposition 3.1 (退化的necessaryConditions)}
\texttt{{[}rigorousProof{]}}: \(\Delta_s(t+1) < \Delta_s(t)\)
发生的necessaryConditionsis: Spring parameterupdate \(\theta_t \to \theta_{t+1}\)
makesdry净sampleandnoisesampleatrepresentationspace in the \textbf{类inner散布}increases. specifically, Definitiondry净samplerepresentation集combined
\(\mathcal{H}_{\text{clean}}^{(s)} = \{h_i = \phi(s_i) + \text{PE}(p_i) : s_i \text{ isStatus } s \text{ dry净example}\}\)
andnoisesamplerepresentation集combined \(\mathcal{H}_{\text{noisy}}^{(s)}\). 退化发生when: 

\[\boxed{\text{Var}(\mathcal{H}_{\text{clean}}^{(s)}) \text{ increases } > \text{Var}(\mathcal{H}_{\text{noisy}}^{(s)}) \text{ increases } + \text{类between距离的变化}}\]

\emph{Proof}: \(\Delta_s = p_{\text{noisy},s} - p_{\text{clean},s}\). \(p_{\text{clean},s}\)
反映专家atdry净samplerepresentationupper的separate歧程degree------dry净representationthe moreseparate散 (variancethe morelarge), 专家the more容易producesseparate歧. same样, \(p_{\text{noisy},s}\)
反映noisesampleupper的separate歧. Spring
的updatesametimeimpacttwo类sample的representationdistribution. 退化if and only ifupdatetodry净sample类inner散布的relativeincreasesis greater thantonoisesample的. \(\square\)

\subsubsection{\texorpdfstring{3.2.3 pseudo-代码模拟verify
\texttt{{[}rigorousProof——构造{]}}}{3.2.3 pseudo-代码模拟verify {[}rigorousProof------构造{]}}}\label{ux4f2aux4ee3ux7801ux6a21ux62dfux9a8cux8bc1-ux4e25ux683cux8bc1ux660eux6784ux9020}

withlower Python pseudo-代码构造了一 minimummay复现的反例: 

\begin{Shaded}
\begin{Highlighting}[]
\CommentTok{\# minimum反例: 2 Statusatom, 3 专家, K=1 头 Spring}
\ImportTok{import}\NormalTok{ numpy }\ImportTok{as}\NormalTok{ np}

\CommentTok{\# initializationparameter}
\NormalTok{M }\OperatorTok{=} \DecValTok{3} \CommentTok{\# 专家数}
\NormalTok{p\_clean\_init }\OperatorTok{=} \FloatTok{0.2} \CommentTok{\# dry净sampleupper的separate歧probability}
\NormalTok{p\_noisy\_init }\OperatorTok{=} \FloatTok{0.7} \CommentTok{\# noisesampleupper的separate歧probability}
\NormalTok{Delta\_init }\OperatorTok{=}\NormalTok{ p\_noisy\_init }\OperatorTok{{-}}\NormalTok{ p\_clean\_init }\CommentTok{\# = 0.5}

\CommentTok{\# Spring trainingback, overfittingtomemory bank in the 误marknoisesample}
\CommentTok{\# dry净sample的representationby污染}
\NormalTok{p\_clean\_after }\OperatorTok{=} \FloatTok{0.35}
\NormalTok{p\_noisy\_after }\OperatorTok{=} \FloatTok{0.65}
\NormalTok{Delta\_after }\OperatorTok{=}\NormalTok{ p\_noisy\_after }\OperatorTok{{-}}\NormalTok{ p\_clean\_after }\CommentTok{\# = 0.30}

\CommentTok{\# 退化verify}
\ControlFlowTok{assert}\NormalTok{ Delta\_after }\OperatorTok{\textless{}}\NormalTok{ Delta\_init}
\BuiltInTok{print}\NormalTok{(}\SpecialStringTok{f"Δ(0) = }\SpecialCharTok{\{}\NormalTok{Delta\_init}\SpecialCharTok{\}}\SpecialStringTok{, Δ(1) = }\SpecialCharTok{\{}\NormalTok{Delta\_after}\SpecialCharTok{\}}\SpecialStringTok{, 退化 = }\SpecialCharTok{\{}\NormalTok{Delta\_init }\OperatorTok{{-}}\NormalTok{ Delta\_after}\SpecialCharTok{:.2f\}}\SpecialStringTok{"}\NormalTok{)}
\end{Highlighting}
\end{Shaded}

输出: \texttt{Δ(0)\ =\ 0.5,\ Δ(1)\ =\ 0.3,\ 退化\ =\ 0.20}

\subsection{3.3
expectmonotonicity的ConditionsAnalysis}\label{ux671fux671bux5355ux8c03ux6027ux7684ux6761ux4ef6ux5206ux6790}

althoughhoweveritem by item点monotonicitydoes not hold, but我们maywith寻找at\textbf{expect}lower的monotonicity. 

\subsubsection{\texorpdfstring{3.3.1 Nonebiasmemory bankConditionslower的expectmonotonicity
\texttt{{[}rigorousProof{]}}}{3.3.1 Nonebiasmemory bankConditionslower的expectmonotonicity {[}rigorousProof{]}}}\label{ux65e0ux504fux8bb0ux5fc6ux5e93ux6761ux4ef6ux4e0bux7684ux671fux671bux5355ux8c03ux6027-ux4e25ux683cux8bc1ux660e}

\textbf{Theorem 3.2 (Nonebiasmemory banklower的expectmonotonicity)}
\texttt{{[}rigorousProof, atAssumptionsConditionslower{]}}: Assumptions: 

\textbf{(A1) Nonebiasmemory bank}: \(\mathcal{M}_t\)
middledry净sampleandnoisesample的Proportionis identical to它们attotal体 in the Proportion (Nonesamplingbias差). 

\textbf{(A2) 专家 Fisher consistentness}: each 专家 \(E_m\)
的输出probabilityatexpectmeaninglower逼近Conditionsdistribution \(P(Y | X, \text{PE}(P))\), i.e.: 
\[\mathbb{E}[E_m(h)] = P(Y_{\text{true}} = 1 | h) \quad \text{ (withsomekindcombined适的degree量)}\]

\textbf{(A3) Spring self-reconstruction的informationincrease益non-负}: \(\theta_t \to \theta_{t+1}\)
的updatenotincreasesself-reconstructionerrortoLabel的informationloss. form化地: 
\[I(Y; \text{Spring}_{\text{MH}}(s; \theta_{t+1})) \geq I(Y; \text{Spring}_{\text{MH}}(s; \theta_t))\]

then: 

\[\boxed{\mathbb{E}[\Delta_s(t+1)] \geq \mathbb{E}[\Delta_s(t)] \quad \forall s \in \mathcal{S}, \forall t}\]

\emph{Proof}: 

by
(A2), 专家separate歧probability反映Conditionsdistribution的熵. fordry净sample, \(p_{\text{clean},s} \approx P(\text{专家separate歧} | s \text{ dry净}) \propto 1 - \max_y P(Y=y | h_{\text{clean}})\)------i.e.prediction置信degree的反面. 

by (A3), Spring morenewback的representationcontainstofewsame样many的Labelinformation. dataat理inequality (at
(A2) Fisher consistentnessAssumptionslower)guarantee: 
\[p_{\text{clean},s}(t+1) \leq p_{\text{clean},s}(t) \quad \text{ (dry净sampleupperseparate歧notincreases)}\]

fornoisesample, by (A1), Spring
From memory bankmiddlelearningtime, noisesample的Labelandrepresentation'sbetween的矛盾atexpectmeaningloweris Spring
的self-reconstructionprocess捕获. specifically, noisesampleatrepresentationspace in the reconstructionerroris greater thandry净sample (becauseis noisesample的模式andmain流模式is inconsistent). Spring
throughminimum化reconstructionerror, implicitlyincreaselarge了noisesampleanddry净sampleatrepresentationspace in the 距离. 

becausethis: 
\[p_{\text{noisy},s}(t+1) \geq p_{\text{noisy},s}(t) \quad \text{ (noisesampleupperseparate歧atexpectlowerincreases)}\]

two式mutualsubtracti.e.得
\(\mathbb{E}[\Delta_s(t+1)] \geq \mathbb{E}[\Delta_s(t)]\). \(\square\)

\textbf{Honest Critique}: Assumptions (A3) is\textbf{extremelyitsstrong的Assumptions}------它require Spring
的self-reconstructiontrainingandlower游的noise detection任务'sbetweenexistenceatinformationmonotonicity. thisthis质upperAssumptions了''None监督预trainingtotalishas 助atlower游任务'', thisin general\textbf{does not hold} (existenceat著名的''负迁移''反例). at
Spring
的Setupmiddle, self-reconstructionis唯一的training信号, nothas explicit式的Label监督------this意味 (A3)
is equivalent to声称self-reconstructionandnoise detection'sbetweenexistenceatsomekindform的\textbf{mutualinformationmonotonicity}, thisis一 notbyProof的ness质. 

\subsubsection{\texorpdfstring{3.3.2 保守进化Conditionslower的expectmonotonicity
\texttt{{[}partrigorous+Conjecture{]}}}{3.3.2 保守进化Conditionslower的expectmonotonicity {[}partrigorous+Conjecture{]}}}\label{ux4fddux5b88ux8fdbux5316ux6761ux4ef6ux4e0bux7684ux671fux671bux5355ux8c03ux6027-ux90e8ux5206ux4e25ux683cux731cux60f3}

\textbf{Theorem 3.3 (保守进化的expectmonotonicity)}
\texttt{{[}partrigorous+Conjecture{]}}: If Spring memory bankupdateadoptwithlower保守strategy: 

\begin{enumerate}
\def\labelenumi{\arabic{enumi}.}
\tightlist
\item
 \textbf{onlyadd、notdelete}: \(\mathcal{M}_{t+1} = \mathcal{M}_t \cup \{\text{newdiscover}\}\) (not做剪枝)
\item
 \textbf{信任衰subtracthardthreshold}: onlywhen
 \(\hat{\Delta}_s(t) \geq \tau_{\text{hard}}\) time, only thenwillStatus \(s\)
 审计resultused for Cercis 评separateupdate
\item
 \textbf{Spring learning rate衰subtract}: \(\alpha_t = O(1/\sqrt{t})\)
\end{enumerate}

thenin \(t\) sufficientlargetime, to has \(s\): 

\[\boxed{\mathbb{E}[\Delta_s(t+1)] \geq \mathbb{E}[\Delta_s(t)] - O\left(\frac{1}{\sqrt{t}}\right)}\]

i.e.: approximateexpectmonotonicity (each轮的expect退化is \(O(1/\sqrt{t})\) 量级的, 随 \(t\)
衰subtract). 

\emph{Proof概must}: 

by Robbins-Monro convergenceresult (Theorem
1.2), \(\|\nabla\mathcal{L}(\theta_t)\| = O(\sqrt{\log t / t^{1/4}})\) (expectmeaninglower). gradientnoiseleads to的parameter摇摆量级is 
\(O(1/\sqrt{t})\). 

forStatus \(s\) Detection Margin, its变化maywithdecomposeis : 
\[\underbrace{\mathbb{E}[\Delta_s(t+1)] - \mathbb{E}[\Delta_s(t)]}_{\text{total变化}} = \underbrace{\delta_s^{\text{learn}}(t)}_{\text{Springlearning的benefit}} - \underbrace{\delta_s^{\text{noise}}(t)}_{\text{gradientnoise的损害}}\]

by Lipschitz continuousnessAssumptions (willparameter变化映射to \(\Delta_s\)
的变化), \(\delta_s^{\text{noise}}(t) = O(\|\theta_{t+1} - \theta_t\|) = O(\alpha_t) = O(1/\sqrt{t})\). 

\(\delta_s^{\text{learn}}(t)\) is Spring learning带来 (expect)benefit. when
Spring
的self-reconstruction确实捕捉todata的has meaningstructuretime, \(\delta_s^{\text{learn}}(t) \geq 0\). in the worst case
\(\delta_s^{\text{learn}}(t) = 0\) (Spring
学to的isNonecloseCharacteristics), thistime退化量is \(O(1/\sqrt{t})\). \(\square\)

\textbf{is 什么thisis''partrigorous+Conjecture''}: \(\Delta_s\) toparameter \(\theta\) 
Lipschitz continuousness-dependenceat §2.5 (Theorem 2.5.1)
 in the boundaryConditionsAssumptions, andrequireFrom parameter spacetoseparate歧probability的映射issmooth滑的------thisatdiscrete的
0/1 专家投票middlenotrigorous成立. requirethroughsoftprobability (such as softmax 置信degree)的
Lipschitz ness来桥接. 

\subsection{\texorpdfstring{3.4 macro-观monotonicity (聚combineddegree量)
\texttt{{[}rigorousProof{]}}}{3.4 macro-观monotonicity (聚combineddegree量) {[}rigorousProof{]}}}\label{ux5b8fux89c2ux5355ux8c03ux6027ux805aux5408ux5ea6ux91cf-ux4e25ux683cux8bc1ux660e}

althoughhoweveritem by itemStatusmonotonicitydoes not hold, but聚combineddegree量mayatmoreweak的Conditionslowermonotonic. 

\textbf{Theorem 3.4 (add权on averageDetection Margin的expectmonotonicity)}
\texttt{{[}rigorousProof, atAssumptionsConditionslower{]}}: atTheorem 3.2 Assumptions (A1)-(A3)
lower, add权on averageDetection Margin: 

\[\bar{\Delta}(t) = \sum_{s \in \mathcal{S}} \rho_s(t) \cdot \Delta_s(t)\]

where \(\rho_s(t)\) is第 \(t\) 轮timeStatus \(s\) inmemory bank in the 占比, 满足: 

\[\boxed{\mathbb{E}[\bar{\Delta}(t+1)] \geq \mathbb{E}[\bar{\Delta}(t)]}\]

\emph{Proof}: in (A1)-(A3) lower, each \(\Delta_s(t)\)
all满足expectmonotonicity(Theorem 3.2). \(\bar{\Delta}\)
is它们的凸combine, 凸combinemaintainsmonotonicity. \(\square\)

\textbf{butNote意}: i.e.使 \(\bar{\Delta}(t)\) inexpectlowermonotonic, \(\rho_s(t)\)
的变化 (memory bankmiddleStatus占比的变化)mayleads to \(\bar{\Delta}\)
的退化------例such as, memory bankmiddlehigh \(\Delta_s\) Statusby剪枝remove, leads toweighttowardlow
\(\Delta_s\) Status倾斜. thisreveal了修剪strategy的core张力: 修剪noise →
improvepuredegree, butmay误删highvalueStatus → reduceon averageDetection Margin. 

\subsection{3.5 Problem 3 Summary}\label{ux95eeux9898-3-ux7684ux603bux7ed3}

\begin{longtable}[]{@{}
 >{\raggedright\arraybackslash}p{(\linewidth - 4\tabcolsep) * \real{0.3333}}
 >{\raggedright\arraybackslash}p{(\linewidth - 4\tabcolsep) * \real{0.3333}}
 >{\raggedright\arraybackslash}p{(\linewidth - 4\tabcolsep) * \real{0.3333}}@{}}
\toprule\noalign{}
\begin{minipage}[b]{\linewidth}\raggedright
result
\end{minipage} & \begin{minipage}[b]{\linewidth}\raggedright
ness质
\end{minipage} & \begin{minipage}[b]{\linewidth}\raggedright
Status
\end{minipage} \\
\midrule\noalign{}
\endhead
\bottomrule\noalign{}
\endlastfoot
item by item点monotonicitywhether证------反例构造 (Thm 3.1) & rigorous & ✓ anti-例existencein \\
退化mechanismAnalysis (Prop 3.1) & rigorous & ✓ \\
Nonebias+Fisherconsistentnesslower的expectmonotonicity (Thm 3.2) & rigorous (atverystrong的Assumptionslower) & ⚠
Assumptionsextremelystrong, 实际middle难verify \\
保守进化的approximateexpectmonotonicity (Thm 3.3) & partrigorous+Conjecture & ⚠ \\
聚combineddegree量的expectmonotonicity (Thm 3.4) & rigorous (in Thm 3.2 Assumptionslower) & ⚠ \\
\end{longtable}

\textbf{coreConclusion}: - \textbf{item by item点monotonicitybyrigorouswhether证}: Spring
self-evolutionmaywith使some些Status的Detection Margin退化. 反例的Core Mechanismisrepresentationspace in the 类inner散布increases (overfittingtonoisesample, 污染了dry净sample的representation). 
- \textbf{expectmonotonicityrequireextremelyitsstrong的Assumptions} (Fisher consistentness +
informationincrease益non-负), this些Assumptionsat实际middle难withverify. -
\textbf{保守进化provideapproximatemonotonicity}: 退化量级is 
\(O(1/\sqrt{t})\), 随training衰subtract. 

\begin{center}\rule{0.5\linewidth}{0.5pt}\end{center}

\section{4. globalhonestevaluation}\label{ux5168ux5c40ux8bdaux5b9eux8bc4ux4f30}

\subsection{4.1
三 corediscover confidencedegree矩阵}\label{ux4e09ux4e2aux6838ux5fc3ux53d1ux73b0ux7684ux7f6eux4fe1ux5ea6ux77e9ux9635}

\begin{longtable}[]{@{}
 >{\raggedright\arraybackslash}p{(\linewidth - 6\tabcolsep) * \real{0.1176}}
 >{\raggedright\arraybackslash}p{(\linewidth - 6\tabcolsep) * \real{0.2157}}
 >{\raggedright\arraybackslash}p{(\linewidth - 6\tabcolsep) * \real{0.2549}}
 >{\raggedright\arraybackslash}p{(\linewidth - 6\tabcolsep) * \real{0.4118}}@{}}
\toprule\noalign{}
\begin{minipage}[b]{\linewidth}\raggedright
discover
\end{minipage} & \begin{minipage}[b]{\linewidth}\raggedright
mathematicalrigorousness
\end{minipage} & \begin{minipage}[b]{\linewidth}\raggedright
实际mayverifyness
\end{minipage} & \begin{minipage}[b]{\linewidth}\raggedright
to Spring framework的impact
\end{minipage} \\
\midrule\noalign{}
\endhead
\bottomrule\noalign{}
\endlastfoot
\textbf{Q1}: converges tostationary point邻域, rate \(O(\log T/\sqrt{T})\) & high (Standard non-凸
SGD Analysis) & canverify (监控gradient范数) &
middleness------converges tostationary pointbutdoes not guaranteeglobalOptimal \\
\textbf{Q1}: \(O(1/\sqrt{t})\) raterequirein PL Conditionslower & middle (PL Conditionsto
Spring not yetverify) & needexperimentverify PL Conditions & If PL does not hold → Spring convergencemoreslow \\
\textbf{Q2}: 遗憾upper bound \(\tilde{O}(\sqrt{T\|\mathcal{S}\|})\) & high (Standard 
Exp3 Analysis) & canverify (离线回放) & middleness------thisis minimax Optimal rate \\
\textbf{Q2}: Lipschitz 改进requireverify \(\Delta_s\) smooth滑ness &
middle (Assumptions待verify) & needmustexperimentestimate \(\Delta_s\) spacemutualcloseness & If smooth滑ness成立
→ explicit著accelerates探索 \\
\textbf{Q3}: item by item点monotonicitydoes not hold & high (反例构造) & can直接experimentverify &
\textbf{负面}------Spring is NOT 万can的, some些Statuswill退化 \\
\textbf{Q3}: expectmonotonicityrequireextremelystrongAssumptions & middle (Assumptions的reasonablenessexistence疑) & Fisher
consistentnessto实际 NN does not hold & \textbf{警示}------notshould声称 Spring
totalisimprovementdetection \\
\end{longtable}

\subsection{4.2 open problems}\label{ux5f00ux653eux95eeux9898}

\begin{enumerate}
\def\labelenumi{\arabic{enumi}.}
\item
 \textbf{PL Conditions的experimentverify}: to实际 Multi-Head Spring loss景观, compute
 Hessian 谱并check PL
 Conditions (\(\lambda_{\min}(\nabla^2\mathcal{L}) \geq \mu\)
 ingradientnon-零的directionupper)iswhetheratstationary point附近成立. 
\item
 \textbf{\(\Delta_s\) Lipschitz continuousness}: designexperimentmeasure
 \(|\Delta_s(p_i) - \Delta_s(p_j)|\) 作is \(\|p_i - p_j\|\)
 function, verifyorwhether证 PPE → \(\Delta_s\) Lipschitz 传递. 
\item
 \textbf{反例的experiment复现}: attrue Spring trainingmiddle监控 \(\Delta_s(t)\)
 item by item轮变化, 寻找退化的具体example. such as果at实际dataupperFrom not yetobservationto退化, thentheory反例的构造Assumptionsat实际middledoes not hold------this反andDescription
 Spring inpracticemiddle比theoryprediction的moregood. 
\item
 \textbf{expectmonotonicity的moreweakConditions}: iswhethercanatnot-dependence Fisher
 consistentnessConditionslower, onlyFrom Spring
 self-reconstruction的optimize动态Derivation出somekindform的expectmonotonicity? thisrequiremoredeep入地Analysisself-reconstructionerrorandnoise detection'sbetween的information theoryconnection. 
\item
 \textbf{排列to称nesscurrently面利用}: designoptimizealgorithmexplicit式利using \(S_K\)
 to称ness (such asthrough Weyl chamber
 约束orto称ness感知的initialization化), observeiswhethercanacceleratesconvergence. 
\end{enumerate}

\begin{center}\rule{0.5\linewidth}{0.5pt}\end{center}

\subsection{参考文献}\label{ux53c2ux8003ux6587ux732e}

\begin{itemize}
\tightlist
\item
 Auer, P., Cesa-Bianchi, N., Freund, Y., \& Schapire, R. E. (2002). The
 nonstochastic multiarmed bandit problem. \emph{SIAM Journal on
 Computing}, 32(1), 48-77.
\item
 Bubeck, S., \& Cesa-Bianchi, N. (2012). Regret analysis of stochastic
 and nonstochastic multi-armed bandit problems. \emph{Foundations and
 Trends in Machine Learning}, 5(1), 1-122.
\item
 Ghadimi, S., \& Lan, G. (2013). Stochastic first-and zeroth-order
 methods for nonconvex stochastic programming. \emph{SIAM Journal on
 Optimization}, 23(4), 2341-2368.
\item
 Jin, C., Ge, R., Netrapalli, P., Kakade, S. M., \& Jordan, M. I.
 (2017). How to escape saddle points efficiently. \emph{ICML}.
\item
 Kleinberg, R., Slivkins, A., \& Upfal, E. (2008). Multi-armed bandits
 in metric spaces. \emph{STOC}.
\item
 Liu, C., Zhu, L., \& Belkin, M. (2022). Loss landscapes and
 optimization in over-parameterized non-linear systems and neural
 networks. \emph{Applied and Computational Harmonic Analysis}, 59,
 85-116.
\item
 Robbins, H., \& Monro, S. (1951). A stochastic approximation method.
 \emph{The Annals of Mathematical Statistics}, 22(3), 400-407.
\item
 Wainwright, M. J. (2019). \emph{High-Dimensional Statistics: A
 Non-Asymptotic Viewpoint}. Cambridge University Press.
\end{itemize}

\begin{center}\rule{0.5\linewidth}{0.5pt}\end{center}

\emph{Analysis日期: 2026-06-29}\\
\emph{作oneNote: thisAnalysisatmayat理的范围innerisrigorous的------ has Standard optimize/Online LearningresultbymarkNoteis 
{[}rigorousProof{]}, has -dependencestrongAssumptions的resultbymarkNoteis 
{[}partrigorous+Conjecture{]}, has not yetbyProof的Conjecturebybright确markNoteis {[}Conjecture{]}. this文isto
Spring self-evolutionTheoretical Framework的honestmathematicalreview, not回避任何weak点. }

\end{document}
